\chapter{Package \textit{Mecanum}}

\section{Téléchargement du package}
L'ensemble des programmes pour utiliser les robots Mecanum se trouvent sur Github. Pour les utiliser, cloner le repository suivant : 

\begin{lstlisting}[style=commandstyle]
git clone https://github.com/watson67/mecanum.git
\end{lstlisting}

\noindent Puis, se placer dans le dossier cloné, et utiliser les commandes suivantes :

\begin{lstlisting}[style=commandstyle]
cd mecanum/src
rm -rf build/install/log/
colcon build --symlink-install --cmake-target clean
colcon build --symlink-install
source install/setup.bash
\end{lstlisting}

\noindent Pour ne pas avoir à sourcer le fichier \textit{setup.bash} à chaque ouverture de terminal, utiliser la commande suivante : 

\begin{lstlisting}[style=commandstyle]
echo "source ~/mecanum/install/setup.bash" >> ~/.bashrc
\end{lstlisting}

\section{Utilisation du package}

\noindent Un tutoriel équivalent au format \textit{markdown} est disponible sur le Github suivant : 
\url{https://github.com/watson67/mecanum/tree/main?tab=readme-ov-file#teleop_mecanum}.

